%
%============= CHAPTER BREAK =================%
%
\chapter{Maximus of Tyre (172--192 CE)}%
\vspace*{3mm}%
\justifying%
%
\noindent{}Maximus of Tyre, or Maximus the Sophist, %
was a prominent orator in the second half of the 2nd century CE. %
The \textit{Suda} places him in the reign of Commodus (180--192 CE). %
Manuscript R of his \textit{Orations} says he delivered them %
in Rome during his first visit to the city %
and places the text between works by two mid-2nd century authors, %
“Alcinuous and Albinus”.\endnote{%
Race 2023:ix--x.} %
In his \textit{Chronicon} (232), %
Eusebius of Caesarea describes Maximus becoming prominent as a contemporary of Arrian of Nicomedia, %
Apollonius the Stoic, and Basilides the Scythian, %
and claims that all of them were sometime educators of Marcus Aurelius.\endnote{%
Race 2023:x.} This last datum need not require a date in the 140s or even 150s CE for Maximusʼs writings, %
since the emperor was well-known for his lifelong pursuit of learning and affinity for a wide array of educators. \\[1mm] %

\noindent{}\textbf{συσπαράσσω} / \textit{susparassō} / “to co-tear asunder” %
is extremely rare, with only eight hits in the entire TLG. %
The LSJ entry for this lemma states that it is a later variant of the Attic %
lemma συσπαράττω / \textit{susparattō}, but the only witnesses to the double-\textit{tau} %
spelling of this lemma are Maximus of Tyre in the late 2nd century CE %
and Byzantine authors of the 12th century CE and later (Chrysoberges, Metochites). %
The first chronologically secure, non-Christian occurrence of either lemma spelling %
is Maximus of Tyre. %
The unprefixed corresponding lemmas were both common in Attic composition: %
σπαράττω / \textit{sparattō} (esp. Aristophanes) and σπαράσσω / \textit{sparassō} (esp. Euripides). %
That is to say, while συσπαράττω / \textit{susparattō} may have \textit{sounded} Attic, %
the συ-prefixed lemma (with either spelling) was actually Hellenistic and Second Sophistic in origin. %

Apparently absent from the major databases of Greek papyri and inscriptions, %
this lemma first appears in Maximusʼ \textit{Dissertationes} 7, %
an oration whose title is, “Which Illnesses are More Grave, %
Those of the Body or Those of the Soul?”\endnote{%
Race 2023.} %
Maximus begins the oration with an analogical frame, %
likening the relationship between the ruler and the body politic to the relationship %
between soul and body (\textit{Diss}. 7.1).\endnote{%
This and the following section summaries are modest re-phrasings of the overview provided by Race 2023:142--143.}  %
Diseases of the soul, in the ruler or anyone, are worse in that they have deleterious %
effects on everyone around them (\textit{Diss}. 7.2). %
A healthy-souled ruler can keep even an unhealthy, endangered city strong, %
as Pericles did for Athens during a plague and the Peloponnesian War (\textit{Diss}. 7.3). %
Pherecydes exemplifies a soul staying healthy even when disease afflicts the body (\textit{Diss}. 7.4). %
The exemplary, healthy soul of Philoctetes, despite his “diseased body”, %
“brings salvation to his comrades” (\textit{Diss}. 7.5).\endnote{%
Race 2023:143.} %

In the immediate context of the relevant quotation, %
Maximus quotes \textit{Od}. 18.74, a suitorʼs taunt about %
Odysseusʼ filthy rags. This taunt foregrounds %
the boxing match that reveals his hidden strength and courage. %
For Maximus, the taunting suitor by contrast typifies cowardice: %

\vfill{}
 \begin{quote}ἡ δὲ δειλὴ ψυχὴ κατορωρυγμένη ἐν σώματι, %
 ὡς ἑρπετὸν νωθὲς εἰς φωλεόν, %
 φιλεῖ τὸν φωλεὸν καὶ οὐδεπώποτε θέλει %
 ἀπαλλαγῆναι αὐτοῦ οὐδὲ ἐξερπύσαι, %
 ἀλλὰ καιομένῳ συγκάεται καὶ σπαραττομένῳ \textbf{συσπαράττεται}, %
 καὶ ἀλγοῦντι τῷ σώματι συναλγεῖ καὶ βοῶντι συμβοᾷ. %
 ὦ πούς, ἀφήσω σε; \\[1mm] %
 \textls[-15]{The cowardly soul, buried down by the body, %
 like a lethargic reptile in a den, %
 loves the den and never ever wishes %
 to be freed from it, nor to creep outside. %
 Instead, as the den is burning, it co-burns, %
 and as the den is torn asunder, it \textbf{is co-torn asunder}, %
 and as the body hurts, it co-hurts, %
 and as it cries out, it co-cries. %
 O foot, shall I release you?} \end{quote} % 
 \vfill{}
 \clearpage

This highly rare and evidently new lemma occurs in both canonical Mark 9:20 and canonical Luke 9:42. %
While surrounding content is attested for Marcionʼs \textit{Evangelion} (*Ev 9:40, 9:44), %
this verse was apparently missing. %
These parallel episodes tell of a boy who is afflicted with a terrible disease, %
a body whose father implores Jesus to heal. %

 \vspace*{2mm} 
  \begin{quote}\textls[-15]{καὶ ἤνεγκαν αὐτὸν πρὸς αὐτόν. %
  καὶ ἰδὼν αὐτὸν τὸ πνεῦμα εὐθὺς \textbf{συν\-εσπάραξεν} αὐτόν, %
  καὶ πεσὼν ἐπὶ τῆς γῆς ἐκυλίετο ἀφρίζων. (Mk 9:20)} \\[0.4mm] %
 And they brought him unto him. %
 And seeing him, the spirit immediately \textbf{co-tore asunder} him, %
 and falling upon the ground he was rolling, foaming.\endnote{%
 This characterization may have been partly inspired by the Euripidean depiction of the Fury-plagued Orestes %
 being thrown to the ground and foaming in \textit{IphTaur} 307--308. %
 On the import of this passage for canonical Acts, see Bilby 2026, chp. 6.} \end{quote} % 
  \vspace*{2mm}
 \begin{quote} ἔτι δὲ προσερχομένου αὐτοῦ ἔρρηξεν αὐτὸν τὸ δαιμόνιον καὶ \textbf{συνεσπάραξεν}· %
  ἐπετίμησεν δὲ ὁ Ἰησοῦς τῷ πνεύματι τῷ ἀκαθάρτῳ καὶ %
  ἰάσατο τὸν παῖδα καὶ ἀπέδωκεν αὐτὸν τῷ πατρὶ αὐτοῦ. (Luke 9:44) \\[0.4mm] %
 Now while he was being brought, the demon broke him and \textbf{co-tore asunder}. %
 But Jesus censured the unclean spirit and %
 healed the child and gave him back to his father. \end{quote} % 
 \vspace*{2mm}

If the story is merely about a demon-possessed child and an individual healing, %
the verbs describing the affliction seem overdone. %
But set against the background of Maximusʼ oration, %
both canonical tales read as political fables of a nationʼs destruction under sick-souled governance, %
and of Jesus as a healthy-souled leader capable of bringing healing. %
The broader context deepens the parallels between these texts. %
In canonical Mark, the ailment and torment of the demon-possessed child sounds %
similar to a reptile writhing on the ground. %
A reptilian redaction unique to canonical Luke also likely owes to Maximusʼ influence. %
In both Gospels (here shared with Marcionʼs \textit{Evangelion}), the following exclamation and rhetorical question %
immediately precedes the miracle: \clearpage %

\begin{quote}ὦ γενεὰ ἄπιστος ἕως πότε ἔσομαι μεθ᾽ ὑμῶν; ἕως πότε ἀνέξομαι ὑμῶν; (*Ev 9:41 // Mk 9:19) \\[0.4mm] %
O faithless generation, how long shall I be with you? How long shall I tolerate you?\end{quote} % 
\vspace*{2mm} 
\begin{quote}ὦ γενεὰ ἄπιστος \textit{καὶ διεστραμμένη}, ἕως πότε ἔσομαι πρὸς ὑμᾶς καὶ ἀνέξομαι ὑμῶν; (Lk 9:41) \\[0.4mm] %
O faithless \textit{and twisted} generation, how long shall I be with you and tolerate you? \end{quote} % 
\vspace*{2mm}

\noindent{}As LSJ notes, διαστρέφω / \textit{diastrephō} can mean “to turn different ways, to twist about, distort”. %
Setting this subtle stylistic choice within its broader, diachronic literary contexts %
suggests it conveys the image of a writhing reptile as well. %

Though Marcionʼs \textit{Evangelion} bears no indication of Maximusʼ influence here, %
its exclamation, “O faithless generation” / ὦ γενεὰ ἄπιστος and rhetorical questions %
likely co-inspired the redactor(s) of canonical Mark and canonical Luke at this point.\endnote{%
My previous research on redactions to canonical Mark and canonical Luke has established %
numerous, distinctive overlaps between the two, lending support to a shared redactor of both canonical texts. %
See Bilby 2020/2025, passim.} %
The supplementation thus made a synthetic association between these prophetic sayings %
of Jesus in Marcion's \textit{Evangelion} and Maximusʼ closing exclamation and rhetorical question: %
ὦ πούς, ἀφήσω σε; / “O foot, shall I release you?” %

In a broader perspective, Maximusʼ evident affinity for συ-prefixed verbs, %
particularly in regard to the dependency of the body on the soul, %
may help account for the gratuitous use of these lemmas (many of them neologisms themselves)\endnote{%
See the penultimate chapter, \textit{infra}.} %
in the canonical New Testament and related early Christian literature %
as products of the Second Sophistic. %